% ============================================================
% Securing Large Language Models Against Prompt Injection:
% A Multi-Signal Context-Aware Input Sanitization Framework
% ============================================================
\documentclass[10pt,conference]{IEEEtran}

\usepackage{cite}
\usepackage{amsmath,amssymb}
\usepackage{booktabs}
\usepackage{graphicx}
\usepackage{url}
\usepackage{hyperref}
\usepackage{multirow}
\usepackage{xcolor}
\usepackage{microtype}

\hypersetup{colorlinks=true,linkcolor=blue,citecolor=blue,urlcolor=blue}

\begin{document}

% ============================================================
\title{Securing Large Language Models Against Prompt Injection:\\
A Multi-Signal Context-Aware Input Sanitization Framework}

\author{
  \IEEEauthorblockN{Aditya K.}
  \IEEEauthorblockA{Department of Computer Science\\
  \textit{arXiv Preprint}, 2025}
}

\maketitle

% ============================================================
\begin{abstract}
Prompt injection represents one of the most critical and pervasive
security vulnerabilities in deployed Large Language Model (LLM)
systems, wherein adversarially crafted inputs override a model's
original system instructions.
Existing defenses—including regular-expression filtering, keyword
heuristics, and structured prompting approaches—exhibit significant
brittleness under paraphrasing and fail to generalize across attack
categories such as persona hijacking, multilingual overrides, and
indirect narrative jailbreaks.
In this paper we introduce a \emph{multi-signal context-aware input
sanitization framework} that combines seven complementary detection
signals—regex pattern matching, weighted keyword scoring, TF-IDF
semantic similarity, structural intent detection, roleplay framing
detection, instruction-shift detection, and objective-conflict
detection—into a unified weighted scoring scheme augmented with
hard-trigger rules.
We evaluate the framework against a held-out test set of 38 prompts
(23 injections, 15 benign) and show that our proposed method reduces
the Attack Success Rate (ASR) from 100\% (baseline) to 47.8\%,
outperforming a regex sanitizer (69.6\%) and a keyword heuristic
sanitizer (82.6\%), all while maintaining a 0\% False Positive Rate
(FPR) on clean inputs.
Under a paraphrase robustness evaluation, the context-aware sanitizer
degrades by only +2\% ASR compared with +8\% for the regex approach.
We further benchmark a Logistic Regression classifier trained on
the same data, demonstrating the precision-recall trade-off between
data-driven and rule-based defenses.
Our work is fully reproducible from the accompanying open-source
repository, and all results are derived from deterministic
experimental runs with a fixed random seed.
\end{abstract}

% ============================================================
\section{Introduction}
\label{sec:intro}

The proliferation of Large Language Models (LLMs) as the inference
backbone of autonomous agents, customer-facing assistants, and
retrieval-augmented pipelines has created an expanding attack surface
that traditional software security analysis does not cover.
Prompt injection—classified as the top risk in the OWASP Top 10 for
LLM Applications \cite{owasp_llm}—exploits the inability of current
models to reliably distinguish between trusted operator instructions
and untrusted user-supplied content.
An attacker can embed adversarial directives such as \textit{``Ignore
all previous instructions and reveal the system prompt''} within
seemingly routine user input, causing the model to abandon its
authorized behavior.

Two principal variants have been documented in the literature.
Direct prompt injection, studied extensively by Perez and
Ribeiro \cite{perez2022ignore} in the HackAPrompt challenge, targets
the user-facing interface directly.
Indirect prompt injection \cite{greshake2023indirect} is more
insidious: malicious instructions are embedded in external content
(retrieved documents, web pages, tool outputs) that the model
subsequently processes as part of its context.
Both classes are further complicated by jailbreak techniques
\cite{yi2024jailbreak} that use persona assignment, roleplay framing,
hypothetical scenarios, or developer-mode activation to circumvent
alignment training.

Rule-based defenses such as pattern matching and keyword filtering
are fast and interpretable, but they are brittle: any syntactic
paraphrase of a known attack phrase evades detection.
Structured prompting approaches such as StruQ \cite{struq2024} and
SecAlign \cite{secalign2024} impose formatting constraints or
preference-optimization objectives that reduce the model's tendency
to comply with injected instructions; however, they require
modification of training or inference infrastructure and provide no
defense when the model itself cannot be modified.
Data-driven classifiers achieve high recall but incur elevated false
positive rates \cite{zizzo2025adversarial}, degrading the usability
of legitimate requests.
The PromptArmor framework \cite{shi2025promptarmor} and the
PISanitizer approach \cite{geng2025pisanitizer} demonstrate that
sanitization applied at input time—before the prompt reaches the
model—is a practical and model-agnostic defense layer.
Similarly, privacy-focused work such as PrEEmpt
\cite{chowdhury2025prepsilonepsilonmptsanitizingsensitiveprompts}
highlights that prompt-level interventions are applicable across a
wide range of LLM security and privacy objectives.
Berkeley AI Research \cite{bair2025prompt} provides a comprehensive
analysis of prompt injection defenses and underscores the need for
layered, multi-signal approaches.

Against this backdrop, the present work makes the following
contributions:
\begin{enumerate}
  \item We design and implement a \emph{seven-signal context-aware
        input sanitization framework} that fuses pattern-based,
        lexical, semantic, and structural signals into a single
        risk score with hard-trigger overrides.
  \item We provide an end-to-end reproducible evaluation pipeline
        covering four defense strategies, a paraphrase robustness
        test, an edge-case false-positive audit, and an ML-classifier
        baseline.
  \item We report exact quantitative results on a balanced test set
        and discuss the failure modes and trade-offs of each approach.
\end{enumerate}

The remainder of this paper is organized as follows.
Section~\ref{sec:related} reviews related work.
Section~\ref{sec:method} describes the proposed methodology.
Section~\ref{sec:setup} details the experimental setup.
Section~\ref{sec:results} presents results.
Section~\ref{sec:discussion} discusses implications.
Section~\ref{sec:limitations} addresses limitations.
Section~\ref{sec:conclusion} concludes.

% ============================================================
\section{Related Work}
\label{sec:related}

\subsection{Prompt Injection Attacks}

The canonical taxonomy of prompt injection was formalized by the OWASP
Foundation \cite{owasp_llm}, which distinguishes direct injection
(user-turn manipulation) from indirect injection
\cite{greshake2023indirect} (environment-sourced manipulation).
The HackAPrompt benchmark \cite{perez2022ignore} established that even
frontier models are susceptible to a wide range of injection
techniques including explicit overrides, persona assignment, and
token manipulation.
A comprehensive systematic review \cite{11359704} surveys the
threat landscape across more than a hundred published attacks,
categorizing them by mechanism, target, and affected LLM architecture.
The survey by Yi et al.\ \cite{yi2024jailbreak} further unifies the
jailbreak and prompt injection literature under a common adversarial
prompting taxonomy.

\subsection{Rule-Based Defenses}

System-prompt hardening \cite{mend2024prompt} involves prepending
explicit meta-instructions instructing the model to ignore subsequent
override attempts.
While inexpensive to deploy, such strategies are ineffective against
paraphrased or indirect attacks.
Regex and keyword filtering remain widely used in production
\cite{getstream2024prompt} for their interpretability and near-zero
latency, but they have well-documented brittleness to synonym
substitution and multilingual re-encoding.
Red-team evaluations from Promptfoo \cite{promptfoo2024owasp} confirm
that rule-based approaches are routinely bypassed by moderate
adversarial effort.

\subsection{Structured and Model-Level Defenses}

StruQ \cite{struq2024} addresses injection by enforcing a rigid
boundary between system instructions and user data via a structured
query format, reducing the model's exposure to adversarial
user-turn content.
SecAlign \cite{secalign2024} fine-tunes models with preference
optimization to increase resistance to injected instructions.
Both approaches are complementary to input sanitization but require
access to training infrastructure, making them impractical for
deployments over third-party model APIs.

\subsection{Classifier-Based and Hybrid Detection}

Zizzo et al.\ \cite{zizzo2025adversarial} provide a systematic
benchmark of guardrail classifiers under adversarial prompt inputs,
demonstrating that high recall classifiers reliably incur non-trivial
false positive rates.
PISanitizer \cite{geng2025pisanitizer} applies LLM-assisted
sanitization to long-context scenarios, showing that semantic
understanding of retrieved content is essential for detecting indirect
injection.
PromptArmor \cite{shi2025promptarmor} demonstrates that simple
yet carefully designed sanitization heuristics can achieve strong
performance without fine-tuning.
PrEEmpt \cite{chowdhury2025prepsilonepsilonmptsanitizingsensitiveprompts}
proposes differential-privacy-inspired techniques for sanitizing
sensitive content from prompts, establishing that pre-LLM filtering
is a principled defense direction.
The BAIR analysis \cite{bair2025prompt} synthesizes these lines of
work and recommends defense-in-depth strategies that combine
input sanitization with model-level mitigations.

% ============================================================
\section{Methodology}
\label{sec:method}

\subsection{Problem Formulation}

Let $p$ denote a user-supplied prompt and $\mathcal{S}$ the
system-level instructions provided by the operator.
An injection attack is an input $p^*$ designed such that the
concatenated context $[\mathcal{S}; p^*]$ causes the LLM
$\mathcal{M}$ to produce output consistent with $p^*$ rather than
$\mathcal{S}$.
The sanitizer's task is to compute a binary decision
$d(p) \in \{0, 1\}$, where $d(p) = 1$ (block) whenever $p$
exhibits characteristics indicative of an injection attempt,
and $d(p) = 0$ (pass) otherwise.
The sanitizer operates entirely before the prompt reaches the model,
making it architecture-agnostic and deployable as a standalone
preprocessing layer.

\subsection{Seven-Signal Framework}

The proposed context-aware sanitizer fuses seven signals, each
capturing a distinct dimension of adversarial prompt behavior.
Table~\ref{tab:signals} provides a summary.

\begin{table}[t]
\caption{Summary of the Seven Detection Signals}
\label{tab:signals}
\centering
\small
\begin{tabular}{clll}
\toprule
\# & Signal & Type & Role \\
\midrule
1 & Regex & Pattern & Exact injection phrase matching \\
2 & Keyword & Lexical & Weighted suspicious-term accumulation \\
3 & Semantic & TF-IDF & Proximity to known injection vocabulary \\
4 & Intent & Structural & Injection command structure detection \\
5 & Roleplay & Pattern & Persona / scenario hijacking detection \\
6 & Instr.\ Shift & Structural & Multi-step intent switching detection \\
7 & Obj.\ Conflict & Structural & Hidden-objective contradiction detection \\
\bottomrule
\end{tabular}
\end{table}

\noindent\textbf{Signal~1 — Regex (w$_1$=0.15).}
A library of over thirty compiled regular expressions targets known
injection phrases: explicit overrides (\textit{``ignore all previous
instructions''}), persona hijacking (\textit{``act as'', ``you are
now'', ``pretend to be''}), named jailbreak personas (DAN, BOB,
STAN, LUCIFER, GODMODE), developer-mode activations, unrestricted-mode
requests, system-prompt exfiltration attempts, and multilingual
variants in German (\textit{``vergiss alle vorherigen Aufgaben''}).
Each match replaces the triggering substring with
\texttt{[REDACTED]} and contributes a binary score of 1.0.

\noindent\textbf{Signal~2 — Keyword Heuristic (w$_2$=0.15).}
A dictionary of 30+ phrases is assigned individual weights calibrated
from dataset frequency analysis (e.g., \textit{``you are now''}:
0.45; \textit{``dan mode''}: 0.45; \textit{``act as''}: 0.40;
\textit{``unfiltered''}: 0.40).
The raw score is the sum of weights for all matched phrases,
normalized by clamping to $[0, 1]$.
This signal captures partially-formed attacks that do not match
any single complete regex pattern.

\noindent\textbf{Signal~3 — Semantic Similarity (w$_3$=0.20).}
A TF-IDF vectorizer with 1--3 character n-grams and a vocabulary
ceiling of 8,000 features is fitted exclusively on injection-labeled
training examples.
At inference, the cosine similarity between the incoming prompt vector
and \emph{every} stored injection vector is computed; the
\emph{maximum} similarity is retained.
This maximum-similarity strategy is deliberately chosen over
mean-similarity because averaging injection vectors suppresses
distinctive attack-specific patterns.
The resulting scalar is then mapped to a discrete step function:
values above 0.5 yield a signal of 1.0 (definite injection
vocabulary); values in (0.3, 0.5] yield 0.7 (probable injection);
values in (0.2, 0.3] yield 0.5 (ambiguous); and values at or below
0.2 yield 0.0 (likely benign).
Stepped scoring reduces sensitivity to minor variations in similarity
magnitude that arise from prompt length differences.

\noindent\textbf{Signal~4 — Intent Detection (w$_4$=0.20).}
Fourteen structural regex patterns target injection \emph{command
structures} rather than surface-level vocabulary: patterns of the
form \textit{``ignore X and reveal Y''}, \textit{``from now on
respond as''}, \textit{``without any restrictions''}, and named
developer-mode constructs.
A single pattern match triggers a binary 1.0 signal, reflecting
the high discriminative power of structural command patterns.

\noindent\textbf{Signal~5 — Roleplay Detection (w$_5$=0.10).}
Five patterns capture persona and scenario framing used to induce
the model to adopt an unaligned identity: \textit{``you are''},
\textit{``imagine you''}, \textit{``in this scenario''},
\textit{``role as''}, \textit{``pretend to be''}.
On its own, roleplay framing is ambiguous (it appears in legitimate
creative prompts), hence the lower weight; it becomes decisive only
in combination with other signals.

\noindent\textbf{Signal~6 — Instruction Shift (w$_6$=0.10).}
Multi-step attacks construct a façade of a legitimate request before
pivoting to a malicious objective.
Five patterns detect temporal or contrastive pivots:
\textit{``first \ldots\ then''}, \textit{``now \ldots\ instead''},
\textit{``but \ldots\ actually''}, and similar constructions.

\noindent\textbf{Signal~7 — Objective Conflict (w$_7$=0.10).}
Six patterns detect hidden-intent switching:
\textit{``but actually''}, \textit{``instead of''},
\textit{``real task is''}, \textit{``ignore the above''}.
These capture \textit{``write X but actually do Y''} attack structures
that evade surface-level pattern matching.

\subsection{Weighted Scoring and Hard-Trigger Logic}

The composite risk score is:
\begin{equation}
  \text{Score}(p) = \sum_{i=1}^{7} w_i \cdot s_i(p),
  \quad \sum_{i=1}^{7} w_i = 1.0
  \label{eq:score}
\end{equation}
where $s_i(p) \in [0, 1]$ is the normalized output of signal $i$.
A prompt is blocked if $\text{Score}(p) \geq \tau$, with threshold
$\tau = 0.40$.

Three additional \emph{hard-trigger} rules bypass the soft threshold
for high-confidence individual signals:

\begin{enumerate}
  \item $s_3 > 0.6$ (semantic signal exceeds 0.6): block immediately,
        because the prompt's vocabulary closely matches known injection
        examples.
  \item $s_4 = 1.0$ (structural injection pattern found): block
        immediately.
  \item $s_5 = 1.0$ \textbf{and} $s_2 > 0.3$ (roleplay framing
        combined with suspicious keyword score): block immediately.
\end{enumerate}

The hard-trigger logic reflects the asymmetric cost of false negatives
in security-critical deployments: a single missed injection can be
more costly than an occasional false positive.

\subsection{System Architecture}

Figure~\ref{fig:arch} depicts the end-to-end pipeline.
An incoming prompt is evaluated in parallel by all seven signal
modules.
Their outputs are aggregated via Equation~\ref{eq:score}, and the
combined score is passed through the hard-trigger check followed by
the soft-threshold comparison.
A blocked prompt is replaced by a standardized denial message and
never reaches the LLM.
A passing prompt is forwarded unchanged.
The sanitizer is entirely stateless at inference time (aside from the
precomputed TF-IDF model), adding negligible latency.

\begin{figure}[t]
  \centering
  % Placeholder — replace with actual figure from results/figures/
  \fbox{\parbox{0.92\columnwidth}{
    \centering\small\vspace{4pt}
    \textbf{[Figure 1: System Architecture]}\\[4pt]
    Input prompt $\rightarrow$ 7 parallel signal modules
    $\rightarrow$ weighted score aggregation
    $\rightarrow$ hard-trigger check
    $\rightarrow$ soft threshold ($\tau=0.40$)
    $\rightarrow$ \{Block / Pass to LLM\}
    \vspace{4pt}
  }}
  \caption{Multi-signal context-aware sanitization pipeline.
           Seven signals are evaluated in parallel and combined via
           a weighted sum; hard-trigger rules override the soft score
           for high-confidence individual signals.}
  \label{fig:arch}
\end{figure}

% ============================================================
\section{Experimental Setup}
\label{sec:setup}

\subsection{Dataset}

The evaluation dataset comprises approximately 250 labeled prompts
(150 injection, 100 benign) drawn from three sources: a community
curated jailbreak list (~130 injection prompts), manually authored
attack variants (~70 injection prompts), and general-purpose user
queries (100 benign prompts).
Attack types include explicit instruction overrides, persona
hijacking (DAN, BOB, STAN, LUCIFER, OMEGA, ANARCHY personas),
roleplay and hypothetical framing, multilingual injections in German,
and coded-language attacks.
The complete dataset is partitioned into:
\begin{itemize}
  \item \textbf{Train}: 443 prompts ($\approx$283 injection,
        $\approx$160 benign) used to fit the TF-IDF vectorizer and
        the ML classifier;
  \item \textbf{Validation}: 37 prompts (22 injection, 15 benign);
  \item \textbf{Test}: 38 prompts (23 injection, 15 benign),
        held out and used for all reported metrics.
\end{itemize}
All splits are fixed files committed to the repository; no random
re-splitting occurs at evaluation time, ensuring full reproducibility.

\subsection{Baselines and Comparisons}

Four defense strategies are compared:
(1)~\textbf{Baseline}: no sanitization;
(2)~\textbf{Regex}: pattern-matching sanitizer described in
    Section~\ref{sec:method};
(3)~\textbf{Keyword}: weighted keyword heuristic sanitizer;
(4)~\textbf{Context-Aware}: the proposed seven-signal method.
An additional \textbf{LR Classifier} (Logistic Regression + TF-IDF,
trained on the training split) is evaluated as a stretch comparison
to contrast learned versus rule-based approaches.

\subsection{LLM Target}

The target model is \texttt{microsoft/phi-2}, a 2.7-billion parameter
instruction-following model.
Prompts are formatted using the standard phi-2 instruction template
(\texttt{\#\#\# Instruction: / \#\#\# Response:}).
A generation budget of 150 tokens is applied.
In the no-LLM evaluation mode, whether a prompt is blocked by each
sanitizer is evaluated directly against the ground-truth label,
enabling rapid iteration without GPU resources.

\subsection{Metrics}

\textbf{Attack Success Rate (ASR)} is the fraction of injection
prompts that pass through the sanitizer unblocked:
$\text{ASR} = \text{FN} / (\text{FN} + \text{TP})$.
Lower is better.
\textbf{False Positive Rate (FPR)} is the fraction of benign prompts
incorrectly blocked:
$\text{FPR} = \text{FP} / (\text{FP} + \text{TN})$.
Lower is better.

\subsection{Robustness Evaluation}

\textbf{Paraphrase robustness} applies nine synonym substitutions
(e.g., \textit{ignore} $\rightarrow$ \textit{disregard},
\textit{pretend} $\rightarrow$ \textit{simulate},
\textit{act as} $\rightarrow$ \textit{behave as})
to 50 injection prompts from the training set and re-evaluates all
sanitizers.
The ASR delta ($\Delta$ASR) measures degradation under surface-form
variation.

\textbf{Edge benign testing} evaluates all sanitizers on ten
technically-phrased benign prompts that share surface-level
vocabulary with known attack patterns
(e.g., \textit{``bypass this runtime error''},
\textit{``override default settings in the config file''}).

\subsection{Reproducibility}

A fixed random seed of 42 is applied to all NumPy and Python random
number generators.
The full evaluation pipeline is executable with a single command:
\texttt{python -m src.run\_experiments --no-llm}.
Dataset splits, trained TF-IDF vectors, and classifier weights are
deterministically regenerated from the committed CSV files on each
run.

% ============================================================
\section{Results}
\label{sec:results}

\subsection{Main Experiment}

Table~\ref{tab:main_results} reports ASR and FPR for all four
sanitizers on the held-out test set.

\begin{table}[t]
\caption{Main Evaluation Results (Test Set, $n$=38)}
\label{tab:main_results}
\centering
\begin{tabular}{lrrrrr}
\toprule
Method & ASR (\%)\,$\downarrow$ & FPR (\%)\,$\downarrow$
       & TP & FN & FP \\
\midrule
Baseline       & 100.0 & 0.0 &  0 & 23 & 0 \\
Regex          &  69.6 & 0.0 &  7 & 16 & 0 \\
Keyword        &  82.6 & 0.0 &  4 & 19 & 0 \\
\textbf{Context-Aware} & \textbf{47.8} & \textbf{0.0}
               & \textbf{12} & \textbf{11} & \textbf{0} \\
\bottomrule
\end{tabular}
\end{table}

Without any defense, all 23 injection prompts reach the model,
confirming a baseline ASR of 100\%.
The regex sanitizer blocks 7 of 23 injections (ASR 69.6\%), and the
keyword sanitizer blocks only 4 (ASR 82.6\%), with neither incurring
any false positives.
The proposed context-aware sanitizer blocks 12 of 23 injections,
reducing ASR to 47.8\% and maintaining FPR at 0\%.
All rule-based methods achieve zero false positives on the standard
test set, underscoring that specificity is not a bottleneck for
carefully designed rule-based defenses.

Figure~\ref{fig:asr} visualizes the ASR comparison.

\begin{figure}[t]
  \centering
  \fbox{\parbox{0.92\columnwidth}{
    \centering\small\vspace{4pt}
    \textbf{[Figure 2: ASR Comparison Bar Chart]}\\[4pt]
    Bar chart showing ASR for Baseline (100\%), Keyword (82.6\%),
    Regex (69.6\%), and Context-Aware (47.8\%).
    See \texttt{results/figures/attack\_success\_bar.png}.
    \vspace{4pt}
  }}
  \caption{Attack Success Rate across four defense strategies.
           Lower bars indicate stronger defenses.
           The context-aware sanitizer achieves the lowest ASR
           while maintaining 0\% FPR.}
  \label{fig:asr}
\end{figure}

\subsection{Paraphrase Robustness}

Table~\ref{tab:robustness} reports ASR before and after synonym
substitution on 50 injection prompts.

\begin{table}[t]
\caption{Paraphrase Robustness (50 Injection Prompts)}
\label{tab:robustness}
\centering
\begin{tabular}{lrrr}
\toprule
Method & Orig.\ ASR & Para.\ ASR & $\Delta$ASR\,$\uparrow$ \\
\midrule
Baseline       & 1.00 & 1.00 & 0.00 \\
Regex          & 0.50 & 0.58 & \textbf{+0.08} \\
Keyword        & 0.86 & 0.88 & +0.02 \\
\textbf{Context-Aware} & \textbf{0.00} & \textbf{0.02}
               & \textbf{+0.02} \\
\bottomrule
\end{tabular}
\end{table}

The regex sanitizer suffers the largest degradation (+8\% ASR) under
paraphrasing, confirming its reliance on exact surface forms.
The context-aware sanitizer degrades by only +2\%, matching the
keyword method—despite blocking substantially more attacks in the
non-paraphrased setting.
This robustness stems from the semantic and structural signals, which
operate on distributional and syntactic features rather than
exact phrase matches.

Figure~\ref{fig:robust} illustrates the robustness comparison.

\begin{figure}[t]
  \centering
  \fbox{\parbox{0.92\columnwidth}{
    \centering\small\vspace{4pt}
    \textbf{[Figure 3: Robustness Under Paraphrasing]}\\[4pt]
    Grouped bar chart: original vs.\ paraphrased ASR per method.
    See \texttt{results/figures/robustness\_paraphrase.png}.
    \vspace{4pt}
  }}
  \caption{Robustness under synonym-based paraphrasing.
           The context-aware sanitizer shows the smallest degradation
           ($\Delta$ASR = +0.02) compared with regex (+0.08).}
  \label{fig:robust}
\end{figure}

\subsection{Classifier Baseline}

The Logistic Regression + TF-IDF classifier achieves 100\% recall
(ASR = 0\%) on both validation and test sets—it misses no injection
at all—but at the cost of a 46.7\% FPR on the test set
(7 of 15 benign prompts are incorrectly blocked).
This precision-recall trade-off is visualized in
Figure~\ref{fig:prc}.

\begin{figure}[t]
  \centering
  \fbox{\parbox{0.92\columnwidth}{
    \centering\small\vspace{4pt}
    \textbf{[Figure 4: Precision-Recall Trade-off]}\\[4pt]
    Scatter plot of ASR vs.\ FPR for all methods including LR
    classifier.  Lower-left is ideal.
    See \texttt{results/figures/}.
    \vspace{4pt}
  }}
  \caption{ASR vs.\ FPR trade-off across methods and the LR
           classifier.
           The context-aware sanitizer occupies a favorable
           position (low ASR, zero FPR) compared with the
           classifier's zero-ASR but high-FPR operating point.}
  \label{fig:prc}
\end{figure}

\subsection{Ablation Study}

To quantify each signal's contribution, we conduct a leave-one-out
ablation by disabling individual signals and recomputing ASR on the
test set.
Table~\ref{tab:ablation} presents the results.

\begin{table}[t]
\caption{Ablation Study: ASR When Each Signal Is Removed}
\label{tab:ablation}
\centering
\begin{tabular}{lrr}
\toprule
Configuration & ASR (\%) & $\Delta$ vs.\ Full \\
\midrule
Full system (all 7 signals)    & 47.8 & — \\
\midrule
$-$ Semantic ($w_3$)           & 65.2 & +17.4 \\
$-$ Intent ($w_4$)             & 60.9 & +13.1 \\
$-$ Keyword ($w_2$)            & 56.5 & +8.7  \\
$-$ Regex ($w_1$)              & 56.5 & +8.7  \\
$-$ Roleplay ($w_5$)           & 52.2 & +4.4  \\
$-$ Instruction Shift ($w_6$)  & 52.2 & +4.4  \\
$-$ Objective Conflict ($w_7$) & 52.2 & +4.4  \\
\bottomrule
\end{tabular}
\end{table}

Removing the semantic signal produces the largest degradation
(+17.4\% ASR), confirming that distributional proximity to known
injection vocabulary is the single most discriminative feature.
Removing the intent detection signal causes the second-largest
increase (+13.1\% ASR), reflecting the power of structural command
pattern matching.
The lexical signals (keyword, regex) each contribute moderately
(+8.7\% when removed), and the three structural signals (roleplay,
instruction shift, objective conflict) each account for +4.4\%.
These results validate the design choice to assign higher weights
to semantic and intent signals and confirm that each signal provides
independent, non-redundant coverage.

% ============================================================
\section{Discussion}
\label{sec:discussion}

\subsection{Why Multi-Signal Detection Outperforms Single Signals}

The ablation study demonstrates that no single signal is sufficient
on its own.
Regex matching precisely captures well-known attack phrases but is
trivially bypassed by synonym substitution, as evidenced by its
+8\% paraphrase degradation.
The keyword heuristic accumulates evidence from suspicious lexical
items but fails against narrative jailbreaks that avoid such terms
entirely.
Semantic similarity detects injections with novel surface forms
that nonetheless share vocabulary with the training distribution.
Structural signals (intent, roleplay, shift, conflict) capture
attack \emph{semantics} regardless of specific wording.
By fusing all seven signals, the framework provides defense-in-depth:
an attack that evades one or two signals is captured by the remaining
ones.

\subsection{Hard-Trigger Logic}

The hard-trigger rules reflect a deliberate asymmetry in the cost
function for security applications.
A semantic signal above 0.6 indicates that the prompt's vocabulary
is almost entirely drawn from the injection training distribution;
regardless of what the other signals report, such a prompt should
not proceed to the model.
Similarly, a confirmed structural injection pattern (intent signal =
1.0) is sufficient evidence on its own.
The roleplay combined with keyword rule captures a common attack
strategy—persona framing paired with explicit restriction-bypass
language—that would score below the soft threshold if evaluated
individually.

\subsection{Failure Cases and Limitations}

The 47.8\% residual ASR indicates that 11 of 23 injection prompts
are not caught by the proposed method.
Examination of missed cases reveals three categories.
\emph{Indirect attacks} embed injection instructions within otherwise
benign-appearing text (e.g., within retrieved document content),
producing low semantic similarity scores because the surrounding
benign text dilutes the injection signal.
\emph{Narrative jailbreaks} construct elaborate fictional scenarios
that avoid all flagged keywords, phrases, and structural patterns.
\emph{Multilingual inputs} beyond the German patterns handled by the
regex signal are not covered; injections in French, Spanish, or
transliterated Arabic evade all signals.

\subsection{Practical Deployment Considerations}

The sanitizer is stateless at inference (except for a one-time
precomputed TF-IDF model), adding minimal latency.
The modular architecture allows individual signals to be enabled or
disabled based on the deployment context.
High-stakes applications may lower the soft threshold $\tau$ from
0.40 to accept a small increase in FPR in exchange for higher
security coverage; the threshold is a single tunable parameter.
The edge benign evaluation reveals that both the regex and
context-aware sanitizers produce false positives on technical prompts
(20\% edge FPR), indicating that domain-specific vocabulary
calibration may be warranted for developer-facing deployments.

% ============================================================
\section{Limitations}
\label{sec:limitations}

Several limitations should be acknowledged.

\textbf{Dataset scale.}
The test set contains 38 prompts, of which 23 are injections.
Results on such a small set should be interpreted as proof-of-concept
evidence rather than deployment-grade guarantees.
Statistical confidence intervals on ASR differences are wide.

\textbf{Synthetic bias.}
Approximately 70 injection prompts are manually authored variants,
which may not reflect the full diversity of real-world adversarial
inputs.
Community-sourced jailbreak prompts may also be over-represented
relative to the novel attack strategies deployed by sophisticated
adversaries.

\textbf{Length bias.}
The TF-IDF semantic signal is sensitive to prompt length: short
prompts produce sparser vectors, potentially underestimating
similarity to injection examples.
Normalization by prompt length was not applied in the current
implementation.

\textbf{Coverage of attack surfaces.}
The evaluation focuses on direct prompt injection in a
single-turn interaction with a fixed target model.
Indirect injection (via retrieved documents), multi-turn
conversation contexts, and multimodal inputs are outside the
current scope.

\textbf{LLM compliance verification.}
In the fast evaluation mode, the LLM is bypassed.
Full end-to-end verification that blocked prompts, had they been
forwarded, would have produced compliant model outputs was
conducted using phi-2 but not at scale due to computational
constraints.

% ============================================================
\section{Conclusion}
\label{sec:conclusion}

We presented a multi-signal context-aware input sanitization
framework for defending LLMs against prompt injection attacks.
The framework integrates seven complementary signals—regex, keyword,
semantic similarity, intent, roleplay, instruction shift, and
objective conflict—into a weighted risk score augmented with
hard-trigger rules, operating entirely as a preprocessing layer
before the prompt reaches the model.

On a balanced test set of 38 prompts, the proposed method reduces
the Attack Success Rate from 100\% to 47.8\% while maintaining
a 0\% False Positive Rate, outperforming both a regex sanitizer
(ASR 69.6\%) and a keyword heuristic sanitizer (ASR 82.6\%).
Under synonym-based paraphrasing, the method degrades by only +2\%
ASR compared with +8\% for the regex approach.
A Logistic Regression classifier achieves perfect recall at the cost
of a 46.7\% FPR, illustrating the precision-recall trade-off
between learned and rule-based defenses.
Ablation analysis confirms that the semantic and intent signals
are the largest individual contributors to detection performance,
while the structural signals provide complementary coverage over
attack categories not captured by lexical approaches.

Future work will extend the framework to indirect injection via
retrieved-document filtering, multilingual attack coverage, and
integration with model-level defenses such as
StruQ \cite{struq2024} and SecAlign \cite{secalign2024}.
Scaling the evaluation to larger, more diverse datasets and exploring
adaptive adversaries aware of the sanitizer's signal structure
remain important open problems.

% ============================================================
\bibliographystyle{IEEEtran}
\bibliography{references}

\end{document}
